\documentclass[12pt,a4paper]{article}

% ============================================================
% PACKAGES
% ============================================================
\usepackage[utf8]{inputenc}
\usepackage[T1]{fontenc}
\usepackage[french]{babel}
\usepackage{geometry}
\usepackage{graphicx}
\graphicspath{{images/}}
\usepackage{xcolor}
\usepackage{listings}
\usepackage{booktabs}
\usepackage{array}
\usepackage{hyperref}
\usepackage{fancyhdr}
\usepackage{titlesec}
\usepackage{enumitem}
\usepackage{mdframed}
\usepackage{tikz}
\usepackage{float}
\usepackage[table]{xcolor}

% ============================================================
% MISE EN PAGE
% ============================================================
\geometry{
  top=2.5cm,
  bottom=2.5cm,
  left=2.5cm,
  right=2.5cm
}

% ============================================================
% COULEURS
% ============================================================
\definecolor{primary}{RGB}{30, 100, 180}
\definecolor{secondary}{RGB}{240, 245, 255}
\definecolor{codebg}{RGB}{245, 245, 245}
\definecolor{codeframe}{RGB}{200, 200, 200}
\definecolor{success}{RGB}{39, 174, 96}
\definecolor{warning}{RGB}{230, 126, 34}
\definecolor{darkgray}{RGB}{60, 60, 60}

% ============================================================
% EN-TETE ET PIED DE PAGE
% ============================================================
\pagestyle{fancy}
\fancyhf{}
\fancyhead[L]{\small\textcolor{primary}{\textbf{Projet M1 GIL -- Gestion de Flotte}}}
\fancyhead[R]{\small\textcolor{darkgray}{Universite de Rouen | 2025--2026}}
\fancyfoot[C]{\thepage}
\renewcommand{\headrulewidth}{0.4pt}

% ============================================================
% STYLE DES TITRES
% ============================================================
\titleformat{\section}
  {\Large\bfseries\color{primary}}
  {\thesection.}{0.8em}{}[\titlerule]

\titleformat{\subsection}
  {\large\bfseries\color{darkgray}}
  {\thesubsection.}{0.6em}{}

\titleformat{\subsubsection}
  {\normalsize\bfseries\color{darkgray}}
  {\thesubsubsection.}{0.5em}{}

% ============================================================
% STYLE DU CODE
% ============================================================
\lstset{
  backgroundcolor=\color{codebg},
  frame=single,
  rulecolor=\color{codeframe},
  basicstyle=\ttfamily\small,
  keywordstyle=\color{primary}\bfseries,
  commentstyle=\color{gray}\itshape,
  stringstyle=\color{success},
  breaklines=true,
  showstringspaces=false,
  numbers=left,
  numberstyle=\tiny\color{gray},
  numbersep=8pt,
  tabsize=2,
  captionpos=b
}

% ============================================================
% BOITE INFO
% ============================================================
\newmdenv[
  backgroundcolor=secondary,
  linecolor=primary,
  linewidth=1.5pt,
  leftline=true,
  rightline=false,
  topline=false,
  bottomline=false,
  innerleftmargin=10pt,
  innerrightmargin=10pt,
  innertopmargin=8pt,
  innerbottommargin=8pt
]{infobox}

% ============================================================
% COMMANDE PLACEHOLDER SCREENSHOT
% ============================================================
\newcommand{\screenshot}[2]{%
  \begin{figure}[H]
    \centering
    \begin{tikzpicture}
      \draw[dashed, color=gray, line width=1pt]
        (0,0) rectangle (14cm, 5cm);
      \node[color=gray] at (7cm, 2.7cm)
        {\large\textbf{[ CAPTURE D'ECRAN ]}};
      \node[color=gray] at (7cm, 2.1cm)
        {\normalsize #1};
    \end{tikzpicture}
    \caption{#2}
  \end{figure}
}

% ============================================================
% DOCUMENT
% ============================================================
\begin{document}

% ------------------------------------------------------------
% PAGE DE TITRE
% ------------------------------------------------------------
\begin{titlepage}
  \centering
  \vspace*{1cm}

  {\Large\textbf{Universite de Rouen Normandie}\\[0.3cm]
  \large Master 1 -- Genie Informatique et Logiciel}

  \vspace{1.5cm}
  \rule{\linewidth}{0.5mm}\\[0.4cm]

  {\Huge\bfseries\color{primary} Projet Gestion de Flotte de vehicules\\[0.3cm]}
  {\LARGE\bfseries Rapport de Semaine 5}\\[0.3cm]
  {\large\color{darkgray} Service Evenements \& Alertes}

  \rule{\linewidth}{0.5mm}\\[1.5cm]

  \begin{minipage}{0.45\textwidth}
    \begin{flushleft}
      \textbf{Etudiant :}\\
      Ahcene Amouchas \\
      Florian Pepin \\[0.5cm]
      \textbf{Encadrants :}\\
      Lydia \& Luc
    \end{flushleft}
  \end{minipage}
  \hfill
  \begin{minipage}{0.45\textwidth}
    \begin{flushright}
      \textbf{Annee universitaire :}\\
      2025 -- 2026\\[0.5cm]
      \textbf{Semaine :}\\
      5 / 7
    \end{flushright}
  \end{minipage}

  \vfill
  {\large\today}
\end{titlepage}

% ------------------------------------------------------------
% TABLE DES MATIERES
% ------------------------------------------------------------
\tableofcontents
\newpage

% ------------------------------------------------------------
% INTRODUCTION
% ------------------------------------------------------------
\section{Introduction et objectifs de la semaine}

\begin{infobox}
\textbf{Objectif de la semaine 5 :} Concevoir et implémenter le service de gestion
des alertes et événements, intégrateur central du système, capable d'agréger les
événements de tous les microservices via Kafka et d'exposer une API REST et GraphQL
pour la consultation et la gestion des alertes en temps réel.
\end{infobox}

\vspace{0.5cm}

La semaine 5 marque une étape charnière dans la construction du système~: le
\textbf{service événements} joue le rôle d'agrégateur central. Contrairement aux
services métier précédents (véhicules, conducteurs, maintenance), ce service ne
produit pas d'événements~: il les \textbf{consomme} depuis tous les autres services
et les transforme en alertes persistées, interrogeables et gérables par les
opérateurs de la flotte.

L'architecture choisie repose sur un découplage total~: le service événements ne
connaît pas les services qui émettent~; il se contente de s'abonner aux topics Kafka
correspondants. Cette approche garantit une extensibilité maximale~: l'ajout d'une
nouvelle source d'alertes (le service localisation, en cours de développement par un
autre membre de l'équipe) ne nécessite aucune modification du code existant.

Les grandes étapes réalisées sont les suivantes~:

\begin{enumerate}[leftmargin=*, label=\textbf{\arabic*.}]
  \item Conception du \textbf{modèle de données} des alertes et des énumérations
        associées (\texttt{AlertType}, \texttt{AlertSeverity}, \texttt{AlertStatus})
  \item Implémentation des \textbf{consumers Kafka} pour les topics
        \texttt{flotte.alertes.events} et \texttt{flotte.conducteurs.events}
  \item Définition du \textbf{contrat d'interface} avec le service localisation
        (topic \texttt{flotte.location.events})
  \item Développement de l'\textbf{API REST} complète avec filtres et gestion du
        cycle de vie des alertes
  \item Branchement complet sur l'\textbf{API GraphQL} de la gateway (résolveurs,
        champs imbriqués Vehicle et Driver)
  \item \textbf{Tests unitaires} (8 scénarios, couverture $>$80\%) et tests
        d'intégration via la collection \textbf{Bruno}
  \item Déploiement validé en \textbf{Docker Compose} et \textbf{Minikube} (Helm)
\end{enumerate}

% ============================================================
% SECTION : ARCHITECTURE DU SERVICE
% ============================================================
\newpage
\section{Architecture du service événements}

\begin{infobox}
\textbf{Principe directeur :} Le service événements est un \textit{aggregator}
\textbf{passif}~: il écoute, transforme et stocke. Il ne modifie jamais l'état
des autres services. Toute alerte créée est immuable dans sa nature~; seul son
statut évolue (ACTIVE $\to$ ACKNOWLEDGED $\to$ RESOLVED).
\end{infobox}

\vspace{0.5cm}

\subsection{Structure des packages}

L'organisation du code suit le même schéma en couches que les autres microservices
du projet~:

\begin{center}
\begin{tabular}{|l|l|}
  \hline
  \rowcolor{primary!10}
  \textbf{Package} & \textbf{Responsabilité} \\
  \hline
  \texttt{config}      & Sécurité Keycloak, création automatique de la base \texttt{events\_db} \\
  \texttt{model}       & Entité JPA \texttt{Alert} et ses trois énumérations \\
  \texttt{dto}         & \texttt{AlertResponse} (record immuable), \texttt{AcknowledgeRequest} \\
  \texttt{repository}  & \texttt{AlertRepository} avec requêtes filtrées par statut/sévérité/type \\
  \texttt{service}     & \texttt{AlertService}~: logique métier, cycle de vie des alertes \\
  \texttt{controller}  & \texttt{AlertController}~: exposition REST avec RBAC \\
  \texttt{events}      & Consumers Kafka et payloads locaux découplés \\
  \hline
\end{tabular}
\end{center}

\subsection{Modèle de données}

L'entité \texttt{Alert} centralise toutes les alertes du système, quelle que soit
leur origine. Les références vers les autres services (\texttt{vehicleId},
\texttt{driverId}) sont des \textbf{références logiques}~: il n'existe aucune clé
étrangère inter-services, conformément au principe d'indépendance des microservices.

\begin{lstlisting}[language=Java, caption=Entité Alert avec cycle de vie géré par JPA]
@Entity
@Table(name = "alerts")
public class Alert {

    @Id
    @GeneratedValue(strategy = GenerationType.UUID)
    private UUID id;

    @Enumerated(EnumType.STRING)
    @Column(nullable = false)
    private AlertType type;        // MAINTENANCE_OVERDUE, LICENSE_EXPIRING...

    @Enumerated(EnumType.STRING)
    @Column(nullable = false)
    private AlertSeverity severity; // INFO, WARNING, HIGH, CRITICAL

    @Enumerated(EnumType.STRING)
    @Column(nullable = false)
    private AlertStatus status = AlertStatus.ACTIVE;

    private UUID vehicleId;  // ref logique - pas de FK inter-service
    private UUID driverId;   // ref logique - pas de FK inter-service

    @Column(nullable = false, columnDefinition = "text")
    private String message;

    @Column(columnDefinition = "text")
    private String metadata; // payload JSON brut pour traçabilité

    private UUID acknowledgedBy;
    private OffsetDateTime acknowledgedAt;
    private OffsetDateTime resolvedAt;

    @PrePersist
    protected void onCreate() { createdAt = updatedAt = OffsetDateTime.now(); }

    @PreUpdate
    protected void onUpdate() { updatedAt = OffsetDateTime.now(); }
}
\end{lstlisting}

\subsection{Types d'alertes supportés}

Six types d'alertes couvrent l'ensemble des cas métier identifiés dans la
spécification du projet~:

\begin{center}
\begin{tabular}{|l|l|l|}
  \hline
  \rowcolor{primary!10}
  \textbf{Type} & \textbf{Source} & \textbf{Sévérité typique} \\
  \hline
  \texttt{MAINTENANCE\_OVERDUE}  & Service Maintenance (Kafka)   & WARNING / INFO \\
  \texttt{LICENSE\_EXPIRING}     & Service Conducteurs (Kafka)   & WARNING \\
  \texttt{LICENSE\_EXPIRING}     & Service Conducteurs (Kafka)   & CRITICAL (si expiré) \\
  \texttt{GEOFENCING\_BREACH}    & Service Localisation (futur)  & HIGH \\
  \texttt{SPEED\_EXCEEDED}       & Service Localisation (futur)  & WARNING à CRITICAL \\
  \texttt{VEHICLE\_IMMOBILIZED}  & Service Localisation (futur)  & WARNING \\
  \texttt{DELIVERY\_TOUR\_DELAYED} & Service Localisation (futur) & INFO \\
  \hline
\end{tabular}
\end{center}

% ============================================================
% SECTION : KAFKA CONSUMERS
% ============================================================
\newpage
\section{Consommation Kafka et traitement des alertes}

\begin{infobox}
\textbf{Choix de désérialisation :} Contrairement aux autres services qui utilisent
le \texttt{JsonDeserializer} typé de Spring Kafka, le service événements utilise
un \texttt{StringDeserializer} brut suivi d'une désérialisation Jackson manuelle.
Ce choix permet de consommer des messages de formats hétérogènes provenant de
services différents (Spring Boot vs Node.js futur), sans dépendance aux en-têtes
de type \texttt{\_\_TypeId\_\_} ajoutés par le \texttt{JsonSerializer}.
\end{infobox}

\vspace{0.5cm}

\subsection{Consumer \texttt{flotte.alertes.events} — Alertes maintenance}

Le service maintenance publie sur ce topic des objets \texttt{AlertEvent} \textbf{sans
enveloppe} \texttt{KafkaEventEnvelope}. Le consumer crée une alerte de type
\texttt{MAINTENANCE\_OVERDUE} et mappe la sévérité textuelle (INFO, WARNING, HIGH,
CRITICAL) vers l'énumération interne.

\begin{lstlisting}[language=Java, caption=Consumer des alertes de maintenance directes]
@KafkaListener(topics = "flotte.alertes.events", groupId = "events-group")
public void handleMaintenanceAlert(String rawMessage) {
    MaintenanceAlertPayload payload =
        objectMapper.readValue(rawMessage, MaintenanceAlertPayload.class);

    alertService.createAlert(
        resolveAlertType(payload.alertType()),
        resolveSeverity(payload.severity()),
        payload.vehicleId(),
        null,
        payload.message(),
        rawMessage // payload brut conservé pour traçabilité
    );
}
\end{lstlisting}

\subsection{Consumer \texttt{flotte.conducteurs.events} — Permis conducteurs}

Le service conducteurs publie des \texttt{KafkaEventEnvelope} dont seuls deux types
d'événements intéressent le service événements~: \texttt{LICENSE\_EXPIRING} et
\texttt{LICENSE\_EXPIRED}. Les autres types (DRIVER\_CREATED, etc.) sont ignorés
sans erreur via un filtre sur le champ \texttt{event\_type}.

\begin{lstlisting}[language=Java, caption=Filtrage des événements conducteurs par type]
@KafkaListener(topics = "flotte.conducteurs.events", groupId = "events-group")
public void handleDriverEvent(String rawMessage) {
    JsonNode root = objectMapper.readTree(rawMessage);
    String eventType = root.path("event_type").asText();

    if (!"LICENSE_EXPIRING".equals(eventType)
            && !"LICENSE_EXPIRED".equals(eventType)) {
        log.debug("Evenement conducteur ignore : {}", eventType);
        return;
    }

    DriverEventPayload payload = objectMapper.treeToValue(
        root.path("payload"), DriverEventPayload.class);

    AlertSeverity severity = "LICENSE_EXPIRED".equals(eventType)
        ? AlertSeverity.CRITICAL : AlertSeverity.WARNING;

    alertService.createAlert(AlertType.LICENSE_EXPIRING, severity,
        null, payload.driverId(), buildLicenseMessage(eventType, payload),
        rawMessage);
}
\end{lstlisting}

\subsection{Contrat d'interface avec le service localisation}

Le service localisation (Node.js, en cours de développement) est intégré via un
consumer prêt à l'emploi sur le topic \texttt{flotte.location.events}. Ce consumer
est déjà déployé mais n'est pas encore sollicité~: il attendra silencieusement
les premiers messages sans erreur ni crash.

Le contrat JSON attendu est documenté directement dans le fichier source
\texttt{LocationEventPayload.java}~:

\begin{lstlisting}[language=Java, caption=Contrat JSON attendu depuis le service localisation]
/**
 * Le location-service publie sur "flotte.location.events" :
 * {
 *   "event_type":     "GEOFENCING_BREACH" | "SPEED_EXCEEDED" | "VEHICLE_IMMOBILIZED",
 *   "vehicle_id":     "uuid",
 *   "driver_id":      "uuid",          // optionnel
 *   "latitude":       48.8566,
 *   "longitude":      2.3522,
 *   "speed_kmh":      120.5,           // pour SPEED_EXCEEDED
 *   "speed_limit_kmh": 90.0,           // pour SPEED_EXCEEDED
 *   "zone_name":      "Zone A",        // pour GEOFENCING_BREACH
 *   "message":        "Description",
 *   "occurred_at":    "2024-01-01T10:00:00Z"
 * }
 */
\end{lstlisting}

La logique de calcul de la sévérité pour les dépassements de vitesse est
automatique~: en dessous de 15 km/h d'excès, l'alerte est WARNING~; entre 15 et
30 km/h, elle passe à HIGH~; au-delà de 30 km/h, elle devient CRITICAL.

% ============================================================
% SECTION : API REST
% ============================================================
\newpage
\section{API REST des alertes}

\begin{infobox}
\textbf{Objectif :} Exposer une API REST claire et filtrée permettant aux opérateurs
de consulter et gérer le cycle de vie des alertes~: de leur création automatique
(par Kafka) jusqu'à leur résolution par un technicien ou un manager.
\end{infobox}

\vspace{0.5cm}

\subsection{Endpoints disponibles}

L'ensemble des routes est regroupé sous le préfixe \texttt{/api/alerts}~:

\begin{center}
\begin{tabular}{|l|l|l|l|}
  \hline
  \rowcolor{primary!10}
  \textbf{Méthode} & \textbf{Endpoint} & \textbf{Description} & \textbf{Rôle requis} \\
  \hline
  GET   & \texttt{/api/alerts}                  & Liste avec filtres optionnels  & Tout utilisateur \\
  GET   & \texttt{/api/alerts/\{id\}}           & Détail d'une alerte            & Tout utilisateur \\
  GET   & \texttt{/api/alerts/stats}            & Nombre d'alertes actives       & Tout utilisateur \\
  PATCH & \texttt{/api/alerts/\{id\}/acknowledge} & Acquitter une alerte         & admin, manager \\
  PATCH & \texttt{/api/alerts/\{id\}/resolve}   & Résoudre une alerte            & admin, manager, technicien \\
  \hline
\end{tabular}
\end{center}

\vspace{0.3cm}

Le endpoint \texttt{GET /api/alerts} supporte les paramètres de filtre suivants~:
\texttt{status}, \texttt{severity}, \texttt{type}, \texttt{vehicleId} et
\texttt{driverId}. Un seul filtre peut être appliqué à la fois~; en l'absence de
filtre, toutes les alertes sont retournées triées par date de création décroissante.

\subsection{Gestion du cycle de vie}

Le cycle de vie d'une alerte suit quatre états, dont trois sont atteignables par
l'API~:

\begin{center}
\begin{tikzpicture}[
  node distance=3.5cm,
  state/.style={draw, rounded corners, fill=secondary, minimum width=2.5cm, minimum height=0.9cm, font=\small\bfseries}
]
  \node[state, fill=warning!20] (active) {ACTIVE};
  \node[state, right of=active] (ack) {ACKNOWLEDGED};
  \node[state, right of=ack, fill=success!20] (resolved) {RESOLVED};
  \node[state, below of=active, node distance=1.8cm, fill=gray!20] (expired) {EXPIRED};

  \draw[->, thick] (active) -- node[above, font=\tiny] {acknowledge} (ack);
  \draw[->, thick] (ack) -- node[above, font=\tiny] {resolve} (resolved);
  \draw[->, thick, bend right=30] (active) to node[right, font=\tiny] {resolve} (resolved);
  \draw[->, dashed] (active) -- node[right, font=\tiny] {auto} (expired);
\end{tikzpicture}
\end{center}

\subsection{Authentification centralisée avec Keycloak}

Une particularité de l'endpoint \texttt{acknowledge} est que l'identifiant de
l'utilisateur qui acquitte n'est \textbf{pas passé dans le corps} de la requête.
Il est extrait directement du jeton JWT via l'annotation Spring Security
\texttt{@AuthenticationPrincipal}. Cette approche évite tout risque d'usurpation
d'identité et simplifie le client (la gateway GraphQL n'a pas à connaître
l'identité de l'utilisateur).

\begin{lstlisting}[language=Java, caption=Extraction du userId depuis le JWT Keycloak]
@PatchMapping("/{id}/acknowledge")
@PreAuthorize("hasAnyRole('admin', 'manager')")
public AlertResponse acknowledge(@PathVariable UUID id,
                                 @AuthenticationPrincipal Jwt jwt) {
    UUID userId = UUID.fromString(jwt.getSubject()); // claim "sub" Keycloak
    return alertService.acknowledge(id, userId);
}
\end{lstlisting}

\subsection{Tests Bruno de l'API REST}

La collection Bruno a été enrichie de neuf fichiers de tests couvrant l'ensemble
des endpoints REST du service, incluant des assertions automatisées sur les codes
de retour et la structure des réponses.

\screenshot{Bruno — Collection service-events : 9 requêtes REST}{Tests Bruno REST du service événements}

\begin{center}
\begin{tabular}{|l|l|l|c|}
  \hline
  \rowcolor{primary!10}
  \textbf{Requête} & \textbf{Endpoint} & \textbf{Code attendu} & \textbf{Résultat} \\
  \hline
  GET   & /api/alerts                         & 200 OK   & \textcolor{green!60!black}{\textbf{OK}} \\
  GET   & /api/alerts?status=ACTIVE           & 200 OK   & \textcolor{green!60!black}{\textbf{OK}} \\
  GET   & /api/alerts?severity=CRITICAL       & 200 OK   & \textcolor{green!60!black}{\textbf{OK}} \\
  GET   & /api/alerts?type=MAINTENANCE\_OVERDUE & 200 OK & \textcolor{green!60!black}{\textbf{OK}} \\
  GET   & /api/alerts/stats                   & 200 OK   & \textcolor{green!60!black}{\textbf{OK}} \\
  GET   & /api/alerts/\{id\}                  & 200 OK   & \textcolor{green!60!black}{\textbf{OK}} \\
  GET   & /api/alerts?vehicleId=\{id\}        & 200 OK   & \textcolor{green!60!black}{\textbf{OK}} \\
  PATCH & /api/alerts/\{id\}/acknowledge      & 200 OK   & \textcolor{green!60!black}{\textbf{OK}} \\
  PATCH & /api/alerts/\{id\}/resolve          & 200 OK   & \textcolor{green!60!black}{\textbf{OK}} \\
  \hline
\end{tabular}
\end{center}

% ============================================================
% SECTION : GRAPHQL
% ============================================================
\newpage
\section{Intégration GraphQL via l'API Gateway}

\begin{infobox}
\textbf{Objectif :} Exposer les alertes dans le schéma GraphQL unifié de la gateway,
en résolvant automatiquement les relations imbriquées \texttt{Vehicle} et
\texttt{Driver} à partir des UUID stockés dans l'alerte.
\end{infobox}

\vspace{0.5cm}

\subsection{Schéma GraphQL}

Le fichier \texttt{alert.graphql} définit le type \texttt{Alert} avec ses champs
imbriqués, deux queries, deux mutations et deux subscriptions (ces dernières
restent en attente de l'implémentation WebSocket prévue en semaine 7)~:

\begin{lstlisting}[language=bash, caption=Schéma GraphQL du domaine alertes]
type Alert {
  id: ID!
  type: AlertType!
  severity: AlertSeverity!
  status: AlertStatus!
  vehicle: Vehicle          # Resolu via vehicle-service
  driver: Driver            # Resolu via driver-service
  message: String!
  acknowledged_by: ID
  acknowledged_at: String
  resolved_at: String
  created_at: String!
}

extend type Query {
  alerts(status: AlertStatus, severity: AlertSeverity,
         vehicle_id: ID, limit: Int = 20, offset: Int = 0): AlertPage!
  alert(id: ID!): Alert
}

extend type Mutation {
  acknowledgeAlert(id: ID!): Alert!
  resolveAlert(id: ID!): Alert!
}
\end{lstlisting}

\subsection{Résolveurs et résolution des types imbriqués}

Le fichier \texttt{alertResolvers.ts} implémente les résolveurs de queries,
mutations et types. La particularité notable est la résolution des champs imbriqués
\texttt{vehicle} et \texttt{driver}~: à partir des UUID stockés dans l'alerte, la
gateway interroge les microservices correspondants en parallèle.

\begin{lstlisting}[language=bash, caption=Résolveurs de types imbriqués (TypeScript)]
Alert: {
  vehicle: async (parent, _, ctx) => {
    if (!parent.vehicle_id) return null;
    try { return await ctx.vehicle.getVehicle(parent.vehicle_id); }
    catch { return null; } // Resilience : l'alerte reste visible meme si
  },                       // le vehicule a ete supprime

  driver: async (parent, _, ctx) => {
    if (!parent.driver_id) return null;
    try { return await ctx.driver.getDriver(parent.driver_id); }
    catch { return null; }
  },
},
\end{lstlisting}

\subsection{Migration depuis les stubs}

Avant cette semaine, les queries \texttt{alerts} et \texttt{alert} ainsi que les
mutations \texttt{acknowledgeAlert} et \texttt{resolveAlert} levaient une erreur
\texttt{NOT\_IMPLEMENTED}. Ils sont désormais branchés sur le service événements
via le \texttt{EventsClient} ajouté au contexte GraphQL.

\begin{center}
\begin{tabular}{|l|l|c|}
  \hline
  \rowcolor{primary!10}
  \textbf{Opération GraphQL} & \textbf{Délégation} & \textbf{Statut} \\
  \hline
  \texttt{alerts(\ldots)}      & GET /api/alerts            & \textcolor{green!60!black}{\textbf{OK}} \\
  \texttt{alert(id)}           & GET /api/alerts/\{id\}     & \textcolor{green!60!black}{\textbf{OK}} \\
  \texttt{acknowledgeAlert(id)} & PATCH /api/alerts/\{id\}/acknowledge & \textcolor{green!60!black}{\textbf{OK}} \\
  \texttt{resolveAlert(id)}    & PATCH /api/alerts/\{id\}/resolve & \textcolor{green!60!black}{\textbf{OK}} \\
  \texttt{Alert.vehicle}       & GET /api/vehicles/\{id\}   & \textcolor{green!60!black}{\textbf{OK}} \\
  \texttt{Alert.driver}        & GET /api/drivers/\{id\}    & \textcolor{green!60!black}{\textbf{OK}} \\
  \texttt{alertCreated}        & WebSocket (Subscription)   & \textcolor{warning}{\textbf{Semaine 7}} \\
  \hline
\end{tabular}
\end{center}

\subsection{Tests Bruno GraphQL}

Six requêtes GraphQL couvrent l'ensemble des opérations du domaine alertes depuis
la gateway, dont une requête exploitant les champs imbriqués \texttt{vehicle} et
\texttt{driver} pour valider la résolution cross-services.

\screenshot{Bruno — Collection GraphQL service-events : queries et mutations}{Tests Bruno GraphQL — résolution des champs imbriqués Vehicle et Driver}

% ============================================================
% SECTION : TESTS
% ============================================================
\newpage
\section{Tests automatisés et couverture de code}

\begin{infobox}
\textbf{Objectif :} Garantir la fiabilité de la logique métier du service via une
suite de tests unitaires isolant la couche \texttt{AlertService}, complétée par un
test d'intégration du contexte Spring Boot.
\end{infobox}

\vspace{0.5cm}

\subsection{Stratégie de tests}

L'approche adoptée suit le même schéma que les semaines précédentes~:

\begin{itemize}
  \item \textbf{Tests unitaires} avec \textbf{JUnit 5} et \textbf{Mockito}~:
        isolation complète de la logique métier du \texttt{AlertService} via
        simulation du \texttt{AlertRepository}.
  \item \textbf{Test de chargement du contexte} Spring Boot avec base de données
        H2 en mémoire, validant la configuration complète du service (Kafka,
        JPA, Security).
\end{itemize}

\subsection{Scénarios de tests unitaires}

Huit scénarios couvrent les cas nominaux et les cas d'erreur de l'\texttt{AlertService}~:

\begin{center}
\begin{tabular}{|l|l|c|}
  \hline
  \rowcolor{primary!10}
  \textbf{Test} & \textbf{Comportement vérifié} & \textbf{Résultat} \\
  \hline
  \texttt{createAlert\_shouldSaveAndReturnAlert}       & Création et persistance d'une alerte  & \textcolor{green!60!black}{\textbf{OK}} \\
  \texttt{createAlert\_shouldSetStatusToActive}        & Statut initial ACTIVE obligatoire     & \textcolor{green!60!black}{\textbf{OK}} \\
  \texttt{findAll\_shouldReturnAllAlerts}              & Retour de la liste complète           & \textcolor{green!60!black}{\textbf{OK}} \\
  \texttt{acknowledge\_shouldUpdateStatusAndUser}      & Passage ACTIVE $\to$ ACKNOWLEDGED      & \textcolor{green!60!black}{\textbf{OK}} \\
  \texttt{resolve\_shouldUpdateStatusAndResolvedAt}    & Passage ACTIVE $\to$ RESOLVED          & \textcolor{green!60!black}{\textbf{OK}} \\
  \texttt{findById\_shouldThrowWhenNotFound}           & Exception sur UUID inconnu            & \textcolor{green!60!black}{\textbf{OK}} \\
  \texttt{acknowledge\_shouldThrowWhenAlertNotFound}   & Exception si alerte introuvable       & \textcolor{green!60!black}{\textbf{OK}} \\
  \texttt{resolve\_shouldBeIdempotentWhenAlreadyResolved} & Idempotence de la résolution        & \textcolor{green!60!black}{\textbf{OK}} \\
  \hline
\end{tabular}
\end{center}

\subsection{Rapport JaCoCo}

\screenshot{Rapport JaCoCo — Service Événements : 9 tests, BUILD SUCCESS}{Couverture de code JaCoCo — Service Evenements}

Les résultats obtenus confirment le respect de l'objectif de couverture fixé à 80\%~:

\begin{itemize}
  \item \textbf{9 tests} exécutés, \textbf{0 échec}, \textbf{0 erreur}.
  \item \textbf{19 classes} analysées (modèle, service, controller, consumers, config).
  \item \textbf{Couverture de la couche Service}~: 100\% des chemins nominaux et des
        cas d'erreur validés.
  \item Les consumers Kafka ne sont pas couverts par les tests unitaires~: leur
        logique de désérialisation nécessiterait des tests d'intégration avec un
        embedded Kafka, prévu en semaine 7.
\end{itemize}

% ============================================================
% SECTION : DÉPLOIEMENT
% ============================================================
\newpage
\section{Déploiement Docker et Kubernetes}

\begin{infobox}
\textbf{Objectif :} Valider le fonctionnement du service événements dans les deux
environnements cibles~: développement local via \textbf{Docker Compose} et
déploiement orchestré via \textbf{Minikube} avec les charts Helm.
\end{infobox}

\vspace{0.5cm}

\subsection{Docker Compose}

Le service est intégré au fichier \texttt{compose.yaml} sous le nom \texttt{events},
accessible sur le port \texttt{8083} de la machine hôte. Il dépend de PostgreSQL
(condition \texttt{service\_healthy}), Kafka et Keycloak.

\begin{lstlisting}[language=bash, caption=Extrait de la configuration Docker Compose]
events:
  build:
    context: ./services/events-service
  container_name: flotte-events
  environment:
    - SPRING_DATASOURCE_URL=jdbc:postgresql://postgres:5432/events_db
    - KAFKA_BROKER=kafka:9092
    - KEYCLOAK_JWK_SET_URI=http://keycloak:8080/auth/realms/...
  ports:
    - "8083:8080"
  depends_on:
    postgres:
      condition: service_healthy
    kafka:
      condition: service_started
\end{lstlisting}

La base de données \texttt{events\_db} est créée automatiquement au démarrage
du service via le \texttt{PostgresDatabaseEnvironmentPostProcessor}, suivant le
même mécanisme que les autres microservices.

\subsection{Déploiement Minikube avec Helm}

Trois modifications ont été nécessaires pour intégrer le service aux charts Helm
existants~:

\subsubsection*{fleet-infra — Création de la base de données}

Le script d'initialisation PostgreSQL a été complété avec la création de la base
\texttt{events\_db}~:

\begin{lstlisting}[language=bash, caption=Ajout de events\_db dans le chart fleet-infra]
primary:
  initdb:
    scripts:
      create-databases.sql: |
        CREATE DATABASE driver_db;
        CREATE DATABASE maintenance_db;
        CREATE DATABASE vehicle_db;
        CREATE DATABASE events_db;    # Ajout
        CREATE DATABASE keycloak_db;
\end{lstlisting}

\subsubsection*{fleet-app — Déclaration du microservice}

L'events-service est ajouté à la map \texttt{services} du chart \texttt{fleet-app}.
Le template Helm génère automatiquement le Deployment, le Service Kubernetes et
les probes de santé (\texttt{startupProbe}, \texttt{readinessProbe},
\texttt{livenessProbe})~:

\begin{lstlisting}[language=bash, caption=Entrée events-service dans fleet-app/values.yaml]
services:
  events:
    name: events-service
    image: events-service:latest
    port: 8080
    replicas: 1
    dbUrlKey: EVENTS_DATASOURCE_URL
\end{lstlisting}

\subsubsection*{Mise à jour du script kube.sh}

Le secret \texttt{EVENTS\_DATASOURCE\_URL}, le build Docker et la vérification
du rollout ont été ajoutés au script de déploiement \texttt{kube.sh}~:

\begin{lstlisting}[language=bash, caption=Ajouts dans kube.sh pour le service événements]
# Secret
kubectl create secret generic db-secrets \
  --from-literal=EVENTS_DATASOURCE_URL=jdbc:postgresql://...5432/events_db \
  ...

# Build dans le daemon Docker de Minikube
docker build -t events-service:latest ./services/events-service/

# Verification du rollout
kubectl rollout status deployment/events-service-deployment \
  -n flotte-namespace --timeout=600s
\end{lstlisting}

\subsection{Vérification en environnement Minikube}

\screenshot{kubectl get pods -n flotte-namespace — events-service en Running}{Pod events-service démarré avec succès dans Minikube}

\screenshot{Ingress flotte.local — route /api/events/ accessible}{Route /api/events/ fonctionnelle via l'Ingress Nginx}

% ============================================================
% SECTION : DÉFIS TECHNIQUES
% ============================================================
\newpage
\section{Défis techniques et résolutions}

\subsection{Hétérogénéité des formats de messages Kafka}

Le principal défi de ce service est la consommation de messages Kafka produits
par des services différents avec des formats différents~: le service maintenance
publie un \texttt{AlertEvent} brut (sans enveloppe), tandis que le service
conducteurs publie un \texttt{KafkaEventEnvelope} standard. Le futur service
localisation en Node.js produira du JSON sans en-têtes de type Java.

\textbf{Solution~:} Utilisation du \texttt{StringDeserializer} pour toutes les
valeurs, suivi d'une désérialisation Jackson manuelle dans chaque consumer. Cette
approche découple totalement le service événements des types Java des autres
services et fonctionne nativement avec tout producteur, quel que soit son langage.

\subsection{Coordination avec le service localisation (développement parallèle)}

Le service localisation est développé en parallèle par un autre membre de l'équipe.
Il était nécessaire de ne pas bloquer le développement du service événements
sur cette dépendance tout en préparant l'intégration future.

\textbf{Solution~:} Le consumer \texttt{LocationEventConsumer} est implémenté et
déployé mais attend simplement les premiers messages. Le contrat d'interface est
formalisé dans le fichier \texttt{LocationEventPayload.java} avec une documentation
complète du format JSON attendu. Quand le service localisation publiera sur
\texttt{flotte.location.events}, le traitement sera automatique sans aucune
modification côté events-service.

\subsection{Extraction du userId pour l'acquittement}

Le schéma GraphQL définit \texttt{acknowledgeAlert(id: ID!): Alert!} sans
paramètre \texttt{userId}. Il était nécessaire d'identifier l'utilisateur qui
acquitte sans modifier le contrat GraphQL public.

\textbf{Solution~:} L'identifiant est extrait du claim \texttt{sub} du jeton JWT
Keycloak directement dans le contrôleur Spring via
\texttt{@AuthenticationPrincipal Jwt}. Le jeton est de toute façon transmis par
la gateway dans l'en-tête \texttt{Authorization}~: aucune information n'est perdue,
et le contrat d'API reste propre.

\subsection{Cohérence des noms entre les manifestes Kubernetes}

Les manifestes Kubernetes bruts (\texttt{event-deployment.yaml},
\texttt{event-service.yaml}) présentaient plusieurs incohérences héritées~:
le \texttt{containerPort} était à 8083 au lieu de 8080, le selector
\texttt{app: event-service} ne matchait pas les labels \texttt{app: event} du
Deployment, et l'Ingress référençait \texttt{event-service} alors que le Service
s'appelait différemment.

\textbf{Solution~:} Réécriture complète des deux fichiers avec des noms harmonisés
(\texttt{events-service} partout), port \texttt{8080} cohérent avec l'application
Spring Boot, et mise à jour de l'Ingress.

\end{document}
